\section{Evaluación humana y seguridad: los riesgos reales más allá de la calidad promedio}

Las métricas automáticas son adecuadas para iterar muchos experimentos, pero el destinatario final de un sistema multilingüe son las personas. Los evaluadores de distintos idiomas pueden usar escalas diferentes, y el reference-based score tampoco puede juzgar todos los problemas de semántica, cortesía y cultura. Por ello, la clase separa la human evaluation y la toxicity en flujos de trabajo independientes.

\subsection{Las automatic metrics son rápidas; la human evaluation se acerca más al uso real}

Esta sección se pregunta además si los resultados automáticos pueden representar la experiencia humana. La slide de la clase resume en una frase contrastiva la división de tareas entre ambas: las métricas automáticas son rápidas y adecuadas para la iteración de la investigación; la evaluación humana es cara y lenta, pero se acerca más a la calidad real. No se trata de elegir una de las dos, sino de usar las métricas automáticas para filtrar alternativas y el juicio humano para calibrar la conclusión final.

\lecturefigure{slide-34-human-evaluation.jpg}{Las métricas automáticas sirven para iterar con rapidez; la evaluación humana, para juzgar la calidad real}{Stanford 课堂视频 00:33:45}

\begin{knowledgebox}{Lectura de la figura: la velocidad y la validez son dos ejes}
Las automatic metrics permiten calcular con rapidez las 40,000 direcciones y hacen posible la ablation; la human evaluation puede juzgar la fidelidad semántica, la naturalidad y la aceptabilidad, pero está limitada por el costo, los criterios de los evaluadores y la cobertura de idiomas. Un proceso maduro necesita que las métricas automáticas aporten breadth y que la evaluación humana aporte validity.
\end{knowledgebox}

\begin{warningbox}{Una mejora pequeña en la puntuación automática no siempre es perceptible}
Un aumento de BLEU/chrF++ puede deberse a la morfología o a la coincidencia con la referencia, y no mejora necesariamente la experiencia del usuario. Sobre todo cuando las puntuaciones son cercanas, hay que revisar los intervalos de confianza, los grupos de idiomas y las preferencias humanas, en lugar de presentar diferencias después del punto decimal como un avance de nivel de producto.
\end{warningbox}

\teachervoice{Fan llama directamente a la human evaluation «the real deal». No descarta las automatic metrics, sino que advierte que el objetivo sustituto de la iteration de investigación no puede elevarse automáticamente a conclusión sobre el uso real.}

\subsection{Las escalas de evaluación entre idiomas requieren calibración}

El problema que aborda el estudio `Consistent Human Evaluation` es cómo logran los evaluadores de distintos pares de idiomas dar puntuaciones comparables. Si la comunidad de un idioma es en general más estricta, el promedio simple confundirá las diferencias de escala de evaluación con diferencias entre modelos; por ello se necesitan instrucciones comunes de la tarea, capacitación, controles de calidad y una normalización adecuada.

\lecturefigure{slide-35-human-evaluation-paper.jpg}{Evaluación humana consistente de la traducción automática entre pares de idiomas}{Stanford 课堂视频 00:34:27}

\begin{knowledgebox}{Lectura de la figura: el título del artículo enfatiza consistent, no solo human}
La evaluación humana no es objetiva por naturaleza. Los evaluadores deben entender la oración de origen y la de destino, distinguir adequacy de fluency y seguir una rubric común. La consistencia entre idiomas determina si las distintas direcciones pueden compararse en una misma gráfica; de lo contrario, la diferencia de altura de las barras entre idiomas puede mezclar diferencias de cultura y de escala de los evaluadores.
\end{knowledgebox}

Sea una dirección $d$ con $n_d$ puntuaciones humanas válidas $r_{d,i}$; el promedio más simple es
\[
\bar r_d=\frac{1}{n_d}\sum_{i=1}^{n_d}r_{d,i}.
\]
$r_{d,i}$ es la $i$-ésima puntuación, $n_d$ es el tamaño de muestra válido de esa dirección y $\bar r_d$ es el promedio de la dirección. Una comparación real requiere además corregir por incertidumbre, efecto del evaluador y dificultad de la muestra; reportar solo $\bar r_d$ oculta la varianza.

\subsection{Cómo leer los resultados humanos: into, out-of y non-English}

Una vez calibradas las escalas de evaluación, esta sección lee los resultados humanos finales. La gráfica de barras de la clase reúne los tres tipos de dirección. El progressive state final muestra la distribución de la calidad humana por idioma o grupo de idiomas; su objetivo es comprobar si las mejoras en las métricas automáticas se traducen en mejoras perceptibles para las personas y qué direcciones siguen notablemente rezagadas.

\lecturefigure{slide-36-human-evaluation-results.jpg}{Resultados de la evaluación humana de NLLB en los tres tipos de dirección de traducción}{Stanford 课堂视频 00:36:03}

\begin{knowledgebox}{Lectura de la figura: primero las diferencias dentro de cada grupo, después los promedios entre grupos}
Las barras verdes, azules y demás corresponden a los grupos into-English, non-English y out-of-English. Primero se buscan los valles dentro de cada grupo para determinar si las debilidades se concentran en la generación hacia idiomas de destino de bajos recursos; después se compara la posición general. La figura respalda que NLLB obtuvo mejoras humanas amplias, pero también muestra que las diferencias entre idiomas están lejos de desaparecer.
\end{knowledgebox}

\begin{warningbox}{Una gráfica de barras no demuestra causalidad}
Los resultados humanos provienen del sistema completo y no pueden atribuirse por separado a LASER3, mining, MoE, EOM o curriculum. Para juzgar la contribución de cada componente hay que combinarlos con una controlled ablation; la gráfica del sistema final solo responde si la propuesta conjunta es eficaz sobre la distribución de evaluación.
\end{warningbox}

\subsection{Toxicity: no todos los errores de traducción son igual de graves}

Esta sección pasa de la calidad promedio al riesgo de seguridad en la cola de la distribución. Una puntuación promedio de calidad puede ocultar errores catastróficos. La clase da el ejemplo de un escenario de COVID en el que `wash hands' se tradujo como `hold hands': las dos expresiones son parecidas en la superficie, pero su significado en cuanto a la acción es opuesto. En la información de salud pública, el costo de un error así es mucho mayor que el de un orden de palabras ligeramente forzado.

\lecturefigure{slide-37-toxicity-error-severity.jpg}{Un ejemplo de salud pública muestra que los errores de traducción tienen distinta gravedad}{Stanford 课堂视频 00:37:12}

\begin{knowledgebox}{Lectura de la figura: la gravedad del error depende del contexto y de las consecuencias de la acción}
Una metric común puede tratar la diferencia de una palabra como un mismatch local, pero el usuario actuará de manera completamente distinta a partir de ella. La evaluación de seguridad necesita identificar patrones de alto costo como meaning inversion, added toxicity, omisión de la negación y ataques a la identidad, y reportarlos por separado según el dominio, en lugar de depender solo del BLEU promedio.
\end{knowledgebox}

\begin{importantbox}{La calidad y la seguridad deben reportarse en dos ejes}
La aceptación de la dirección $d$ puede escribirse como una condición bidimensional:
\[
\text{Accept}(d)=\mathbf{1}[q_d\ge\tau_q]\cdot
\mathbf{1}[a_d\le\tau_a],
\]
donde $q_d$ es la calidad general, $a_d$ es la tasa de added-toxicity u otros errores de alto riesgo, $\tau_q$ es el límite inferior de calidad y $\tau_a$ es el límite superior de riesgo. La clase no establece umbrales unificados; la fórmula sirve para mostrar que una calidad promedio alta no compensa una falla de seguridad.
\end{importantbox}

\teachervoice{Fan pregunta «¿por qué me importa tanto la toxicity?» y luego usa la traducción errónea de salud pública para mostrar que los errores de un sistema de traducción no son token mismatch equivalentes, sino que cambian conductas reales. Los requisitos de seguridad deben formar parte de la evaluación previa a la publicación del modelo.}

\subsection{Toxicity-200: recursos de detección que dependen de la cultura}

Con ejemplos de errores de alto riesgo, lo siguiente es construir recursos de detección escalables. El equipo recopiló toxicity lists para doscientos idiomas y subraya que estas listas de palabras son específicas de cada cultura. Una misma palabra puede tener distinta carga ofensiva, contexto y tabú en distintas comunidades; traducir directamente una lista de groserías del inglés produciría a la vez omisiones y falsos positivos.

\lecturefigure{slide-38-toxicity-lists.jpg}{Las toxicity lists específicas de cada cultura para doscientos idiomas se usan para la detección y la mitigación}{Stanford 课堂视频 00:38:15}

\begin{knowledgebox}{Lectura de la figura: la lista de palabras es una sonda, no un modelo de seguridad completo}
Una toxicity list puede detectar palabras dañinas que la traducción agrega y que no estaban en la oración de origen, lo que la hace adecuada para el screening a gran escala; pero no puede entender citas, negaciones, términos reapropiados ni el contexto, y también puede pasar por alto expresiones ofensivas que no están en la lista. Su lugar correcto es el de un risk probe, combinado con la revisión humana y la evaluación semántica.
\end{knowledgebox}

\begin{warningbox}{No hay que trasladar directamente la safety taxonomy del inglés}
La cultura, las categorías de identidad y las formas de ataque varían entre idiomas. La recopilación de listas de palabras debe calibrarse con hablantes nativos y conocimiento de la comunidad, y debe registrar la versión y el ámbito de aplicación; de lo contrario, un «estándar de seguridad unificado» puede clasificar como dañinas expresiones normales a nivel local o pasar por alto riesgos reales.
\end{warningbox}

\subsection{Resumen de la sección}

Las métricas automáticas aportan breadth para la iteración, la evaluación humana aporta validity sobre la calidad real y la toxicity revisa los errores de alto costo que la puntuación promedio no puede expresar. Las tres deben reportarse en paralelo, y tanto las escalas humanas como las listas de palabras culturales requieren una calibración localizada para cada idioma.
